\chapter{Introduzione}

\section{Statement}
Il progetto \textbf{EscapeManager} è un sistema gestionale pensato per una catena in \textit{franchising} di Escape Room. L'idea è centralizzare in un unico software la prenotazione delle stanze da parte dei clienti e il controllo delle sessioni di gioco da parte del personale.

Il dominio dell'applicazione si basa su entità principali come \textbf{Sede}, \textbf{Stanza}, \textbf{Utente} (con vari ruoli) e \textbf{Prenotazione}. Ogni sede ospita più stanze, ognuna con un tema, una difficoltà, una capienza e un ciclo di vita operativo preciso (es. \textit{Disponibile, In Corso, In Pulizia, In Manutenzione}).

Il sistema prevede tre tipi di utenti:
\begin{itemize}
    \item \textbf{Cliente:} può sfogliare il catalogo delle stanze, prenotare una sessione e mettersi in lista d'attesa per uno slot momentaneamente occupato.
    \item \textbf{Game Master (Operatore):} gestisce le operazioni quotidiane. Controlla le prenotazioni, avvia e termina le sessioni, aggiorna lo stato delle stanze (ad esempio bloccandole per manutenzione) e registra l'esito finale della partita.
    \item \textbf{Admin (Amministratore):} ha una visione globale sul franchising e può configurare le politiche di prezzo (ad esempio impostando i prezzi del weekend).
\end{itemize}

Ho organizzato il codice secondo il pattern BCE (Boundary-Control-Entity), applicando i principi dell'Object-Oriented Design. Per i problemi centrali del dominio — come il cambio di stato delle stanze o il calcolo del prezzo dinamico — ho applicato diversi Design Pattern GoF. I dati sono salvati su un database relazionale PostgreSQL.

\section{Stack Tecnologico e Strumenti}
Nello sviluppo sono state utilizzate le seguenti tecnologie:
\begin{itemize}
    \item \textbf{Linguaggio base:} Java (JDK 17).
    \item \textbf{Database:} PostgreSQL. Il codice Java si collega tramite JDBC a un layer DAO per mappare gli oggetti sulle tabelle.
    \item \textbf{Build Tool:} Maven, per gestire le dipendenze e il ciclo di build.
    \item \textbf{Mockup UI:} JavaFX (FXML), usato solo per creare un prototipo visivo slegato dalla business logic.
    \item \textbf{Testing:} JUnit 5 (Jupiter) per i test unitari e di integrazione, e \textbf{Mockito} per simulare il database nei Mock test.
    \item \textbf{Versioning e UML:} Git tramite GitHub, IDE Visual Studio Code, e PlantUML integrato per generare i diagrammi direttamente da formato testo.
\end{itemize}

\section{Note sull'uso dell'Intelligenza Artificiale}
Come previsto dalle direttive del corso, nello sviluppo di \textit{EscapeManager} ho usato strumenti di AI generativa (GitHub Copilot e Google Gemini) come supporto alla programmazione. 

Non ho delegato il lavoro all'AI, ma l'ho usata come supporto pratico:
\begin{itemize}
    \item \textbf{Decisioni architetturali:} Le scelte di design (uso del pattern BCE, o dove applicare State e Strategy) derivano dallo studio teorico fatto a lezione.
    \item \textbf{Scrittura codice:} Ho chiesto all'AI di aiutarmi ad accelerare la scrittura di codice ripetitivo (ad esempio i metodi getter/setter, l'ossatura delle classi DAO o la base per i test JUnit).
    \item \textbf{Revisione testo e \LaTeX:} L'AI è stata utile per impaginare rapidamente i diagrammi e il formattare il testo in \LaTeX.
    \item In ogni caso, ogni frammento di codice o testo suggerito è stato letto, compreso e spesso modificato prima di essere inserito nel progetto.
\end{itemize}

\section{Struttura del documento}
Nei capitoli che seguono: il 2 copre l'analisi dei requisiti con Use Case e mockup; il 3 presenta le scelte di progettazione e i diagrammi UML; il 4 mostra i frammenti di codice Java più significativi; il 5 descrive la strategia di test.

\section{Stato del prototipo e limitazioni}
\textit{EscapeManager} è un prototipo didattico funzionante, focalizzato sulla logica di business e sulla persistenza dei dati. L'interazione avviene tramite un menu a linea di comando (CLI); le schermate grafiche presenti nella relazione sono mockup che illustrano l'interfaccia ipotizzata.

L'infrastruttura di base (come l'iscrizione a una lista d'attesa modellata con il pattern Observer) è testata e implementata correttamente lato codice. Tuttavia, alcune dinamiche reali (come l'invio concreto di un'email tramite SMTP o un gateway di pagamento) esulano dal perimetro prototipale in esame e non sono state implementate.
